\documentclass[12pt, titlepage]{article}
%\usepackage[norsk]{babel}
\usepackage[utf8]{inputenc}
\usepackage{minted}
\usepackage[colorlinks = true,
linkcolor = blue,
urlcolor = blue]{hyperref}
\usepackage[norsk]{babel}
\usepackage[margin=1in]{geometry}
\usepackage{array}
\usepackage{booktabs}
\usepackage{pdflscape}

\begin{document}

\title{Sensurveiledning}
\author{IIRA2001 - Høst 2024}
\date{}
\maketitle

Til sensuren har vi tatt utgangspunkt i Karakterbeskrivelsen til NTNU for teknologiske fag (Tabell \ref{tab:karakterbeskrivelse})

\begin{table}[h!]
   \footnotesize
   \begin{center}
   \begin{tabular}{|l|l|p{12cm}|}
      \hline
      \textbf{Symbol}   &  \textbf{Betegnelse}  &   \textbf{Utdypende beskrivelser av vurderingskriterier} \\
      \hline
      A & Fremragende & Kandidaten viser særdeles god kunnskap i og oversikt over emnets faglige grunnlag og innhold. Kandidaten viser meget stor grad av selvstendighet og solid analytisk forståelse. Kandidaten viser svært gode ferdigheter i anvendelsen av denne kunnskapen\\
      \hline
      B & Meget god   &Kandidaten viser meget god kunnskap i og oversikt over emnets faglige grunnlag og innhold. Kandidaten viser betydelig grad av selvstendighet og god analytisk forståelse. Kandidatens ferdigheter i anvendelsen av denne kunnskapen ligger over gjennomsnittet.\\
      \hline
      C & God &Kandidaten viser god oversikt over de viktigste kunnskapselementene og sammenhengene i emnets faglig grunnlag og innhold. Kandidaten viser selvstendighet. Kandidaten viser analytisk evne og forståelse. Kandidaten viser gjennomsnittlig evne til å anvende sin kunnskap. Gjennomsnittet kan dels tolkes som typisk prestasjon av mange studenter i emnet og dels som krav til tilfredsstillende god prestasjon i emnet.\\
      \hline
      D & Nokså God &Kandidaten viser i noen grad analytisk evne og forståelse. Kandidaten viser selvstendighet i noen grad. Kandidaten viser oversikt over de viktigste kunnskapselementene og sammenhengene i emnets faglige innhold, men denne oversikten er preget av noen vesentlige mangler. Kandidaten viser i noen grad evne til å bruke kunnskapen aktivt, men prestasjonen er noe dårligere enn gjennomsnittet.\\
      \hline
      E & Tilstrekkelig & Kandidaten viser mangelfull analytisk evne og forståelse. Kandidaten viser gjennomgående noe, men sporadisk preget oversikt over de viktigste kunnskapselementene og sammenhengene i emnets faglige innhold. Kandidatens prestasjon oppfyller minimumskravet som stilles i emnet når det gjelder kunnskap, analytisk evne og ferdighet i å anvende emnets kunnskapsinnhold.\\
      \hline
      F & Ikke bestått & Kandidatens prestasjon faller under minimumskravet som stilles i emnet når det gjelder kunnskap, analytisk evne og ferdighet i å anvende emnets kunnskapsinnhold. \\
      \hline
   \end{tabular}
   \end{center}
   \caption{Karakterbeskrivelse for emner i teknologiske fag}\label{tab:karakterbeskrivelse}
\end{table}

\newpage
\subsection*{Sensurveiledning}
\subsubsection*{F}
Kandidaten har ikke tilegnet seg nok programmeringsferdigheter til å lage programmer i python som gjør enkle beregninger.

\subsubsection*{E}
Kandidaten tilfredsstiller minimumskravet. Kandidaten har klart å skrive kode som gjør enkle beregninger i programmene som omhandlet populasjonvekst og sparing.
Kandidaten viser ikke vesentlige ferdigheter innen dataanalyse med pandas eller bruk av web-apier.

\subsubsection*{D}
Kandidaten viser at hen kan skrive kode som gjør enkle beregninger i programmene som omhandlet populasjonvekst og sparing, men sliter gjør det ikke helt rett når det kommer til simulering eller dataanalyse. 
Kanditaten viser en del kompetanse rundt mange av de generelle konseptene som er gjennomgått, men har vesentlige mangler når det kommer
til dataanalyse med pandas eller bruk av web-apier.

\subsubsection*{C}
Kandidaten viser gode programmeringsferdigheter og har flittig brukt mange av de generelle programmeringskonseptene vi har jobbet med. 
Kandidaten behersker også plotting med pyplot og enkel bruk av de mer avanserte dataanalyseverktøyene (pandas) og/eller web-apier. Kandidaten er i stand
til å bruke programmering noe selvstendig til et eget faglig problem, og kan tildels drøfte og analysere resultatene de får fra programmene sine.

\subsubsection*{B}
Kandidaten har meget gode programmeringsferdigheter og viser god kontroll over både de generelle konseptene og de mer avanserte dataanalyseverktøyene.
Kandidaten viser å kunne bruke python selvstendig som et verktøy i en meningsfull faglig sammenheng, og kan drøfte og analysere resultatene fra programmene sine.
Kandidaten fremstiller informasjon og data godt både grafisk og tekstlig.

\subsubsection*{A}
Kanidaten viser fremragende programmeringsferdigheter og meget god bruk av dataanalyseverktøy og/eller web-apier. 
Kandidaten har også vist meget høy grad av selvstendighet og har nyttiggjort programmering i en noe utfordrene faglig sammenheng.
Kandidaten drøfter og analyerer resultatene fra programmene meget godt.

\subsection*{Merknader}
Oppgaven til siste del av mappen var åpen for tolkning og valg av fremgangsmåte. Vi vil derfor prøve å være noe fleksible i sensurkravene over. 
Vi har lagt en del vekt på at kandidaten viser bruk av programmering som verktøy i en faglig sammenheng og drøfter og analyserer resultatene, spesielt
for å oppnå en 'A' eller 'B'. Avhengig av hvilken vinkling man har hatt på oppgavene sine faller dette mer eller mindre naturlig. Om besvarelsen for eksempel er
av en mer teknisk karakter med fokus på programmeringen - eller er med praktisk rettet uten et spesifikt problem å anvendes på - passer kanskje ikke vurderingskriteriene over godt nok.
Vi skal prøve å ta hensyn til dette.

Det er også mulig å vise høy grad av selvstendighet i en faglig relevant sammenheng og drøfte og analysere programmet sitt godt, men falle litt kort når det kommer til programmeringsferdihetene.
Vi vil i slike tilfeller ta en helhetlig vurdering og bruke refleksjonotatet, til å komme frem til endelig sensurvurdering.

\section*{Programmeringsferdigheter}
Når vi vurderer programmeringsferdigheter ser vi spesielt etter hvordan kandidaten bruker variabler, funksjoner, datastrukturer og iterasjon.
Følgende rubrikk anvendes ikke direkte under sensuren, men kan være en god veiledning for hvordan vurdering foregår.
\begin{landscape}

\tiny
\begin{table}[h!]
\centering
\begin{tabular}{p{4cm} | p{4cm} | p{4cm} | p{4cm} | p{4cm}}
\toprule
\textbf{Kriterier} & \textbf{Lav oppnåelse (1)} & \textbf{Under utvikling (2)} & \textbf{God (3)} & \textbf{Avansert (4)} \\
\midrule

\textbf{Bruk av variabler} & 
Bruker variabler feil eller ikke i det hele tatt; variabelnavn er ikke beskrivende eller forvirrende. &
Bruker variabler, men kan ha inkonsekvente eller uklare navn; kan misbruke variabler. &
Bruker konsekvent variabler riktig; variabelnavn er beskrivende og følger konvensjoner. &
Demonstrerer avansert bruk av variabler; anvender beste praksis for variabelscope, konstanter og navnekonvensjoner; koden er klar og vedlikeholdbar. \\

\midrule

\textbf{Bruk av funksjoner} &
Bruker ikke funksjoner eller misbruker innebygde funksjoner; koden er repetitiv. &
Bruker innebygde funksjoner, men definerer sjelden egne funksjoner; kan ha problemer med parametere eller returverdier. &
Definerer og bruker egne funksjoner; forstår parametere, argumenter og returverdier&
Skaper effektive og gjenbrukbare funksjoner; Bruker funksjoner til å lage en god programstruktur \\

\midrule

\textbf{Håndtering av datastrukturer} &
Viser begrenset forståelse av datastrukturer; kan misbruke eller unngå dem. &
   Bruker grunnleggende datastrukturer (lister og dictionaries) &
Velger og bruker passende datastrukturer effektivt; demonstrerer forståelse for deres egenskaper og metoder. &
Viser avansert bruk av datastrukturer som dictionaries med komplisert struktur; Kan tilpasse bruken etter behov; Bruker listekomprehensjon fornuftig \\

\midrule

\textbf{Håndtering av iterasjon} &
Sliter med iterasjon; kan ikke bruke løkker riktig eller i det hele tatt. &
Bruker løkker, men kan ha ineffektiviteter eller feil;&
Bruker løkker og iterasjon effektivt; unngår vanlige feil. &
Demonstrerer avanserte iterasjonsteknikker; bruker list comprehension, og anvender iterasjon for å behandle data effektivt. \\

\bottomrule
\end{tabular}
   \caption{Typisk rubrikk for vurdering av generelle programmeringsferdigheter (merk overlapp i datastrukturer og iterasjon)}
\end{table}


\end{landscape}
\end{document}





\end{document}
